\documentclass[11pt]{article}
\usepackage[a4paper,margin=22mm]{geometry}
\usepackage{fontspec}
\setmainfont{DejaVu Serif}
\setsansfont{DejaVu Sans}
\usepackage{microtype}
\usepackage{xcolor}
\usepackage{hyperref}
\usepackage{enumitem}
\usepackage{parskip}
\definecolor{Accent}{HTML}{971D32}
\definecolor{Muted}{HTML}{666666}
\hypersetup{colorlinks=true,urlcolor=Accent,linkcolor=Accent}
\setlist{leftmargin=1.6em,itemsep=0.3em,topsep=0.35em}
\setcounter{tocdepth}{2}
\renewcommand{\contentsname}{Contents}
\newcommand{\sectionrule}{\par\vspace{-0.5em}\noindent\textcolor{Accent}{\rule{\linewidth}{0.6pt}}\par}
\begin{document}

\begin{center}
{\sffamily\bfseries\color{Accent} TOEFL WORD OF THE DAY\par}
\vspace{8mm}
{\fontsize{42}{48}\selectfont\bfseries trite\par}
\vspace{2mm}
{\Large\sffamily\color{Accent} /traɪt/\par}
\vspace{2mm}
{\small\sffamily Local audio phrase: trite\par}
{\small\sffamily Audio file: audios/trite.mp3\par}
\vspace{2mm}
{\large\sffamily\color{Muted} adjective\par}
\vspace{5mm}
{\sffamily overused and no longer original\par}
\vspace{6mm}
{\large Trite describes words, ideas, or expressions that have been repeated so often that they have lost their freshness or impact.\par}
\vspace{6mm}
\begin{minipage}{0.84\textwidth}
\itshape “His advice to simply follow your dreams sounded trite to the students.”
\end{minipage}\par
\vspace{10mm}
\textbf{Date:} August 20th 2026\quad\textbullet\quad \textbf{Level:} B1--mid-B2 TOEFL
\vfill
Created and compiled by \textbf{Amir Kooshky}\par
\vspace{2mm}
\href{https://t.me/KooshkyTOEFL}{Telegram Channel}\quad\textbullet\quad
\href{https://www.instagram.com/kooshkytoefl}{Instagram}\quad\textbullet\quad
\href{https://t.me/KooshkyTOEFL\_pv}{Personal Telegram}
\end{center}
\newpage
\tableofcontents
\newpage

\section{Meanings and usage}
\sectionrule
\subsection{Meaning 1: Overused and no longer original}
\textbf{Part of speech:} adjective

\textbf{IPA:} /traɪt/

\textbf{Pronunciation phrase:} trite

\textbf{Local audio file:} audios/trite.mp3

\textbf{Register:} neutral to formal; usually disapproving and common in discussions of writing, speech, art, advertising, and ideas

\textbf{Definition:} expressed or used so often that it no longer seems interesting, original, or effective

\textbf{In simpler words:} used so much that it feels boring or unoriginal

\textbf{Typical pattern:} trite + phrase, expression, saying, remark, observation, idea, or argument

\subsubsection{Natural examples}
\begin{enumerate}
  \item The essay relies on trite phrases instead of developing its ideas clearly.
  \item His advice to simply follow your dreams sounded trite to the students.
  \item The advertisement was criticized for using trite images of success and happiness.
  \item The speaker avoided trite statements and supported her argument with specific evidence.
  \item Although the message was sincere, its wording seemed trite and predictable.
  \item The novel contains several trite observations about love and relationships.
  \item Teachers encouraged students to replace trite expressions with more precise language.
  \item The conclusion that technology has both advantages and disadvantages can sound trite unless it is developed with specific examples.
\end{enumerate}

\subsubsection{Nuance and implication}
\textbf{Central nuance:} Trite does not simply mean common or simple. It means something has been repeated so often that it has lost much of its originality, freshness, or emotional effect.

\textbf{Usual implication:} The speaker or writer is criticizing the expression or idea as predictable and insufficiently thoughtful.

\textbf{Compared with original:} An original expression or idea feels fresh or distinctive. A trite one feels familiar because it has been used too many times.

\subsubsection{Grammar notes}
\textbf{How to use it:} Use it before a noun, as in a trite expression, or after a linking verb, as in The conclusion sounds trite.

\textbf{What can follow it:} It commonly describes phrases, expressions, sayings, remarks, observations, ideas, arguments, themes, and advice.

\textbf{Position in a sentence:} It often appears after sound, seem, become, or feel when judging language or an idea.

\subsubsection{Useful patterns}
\textbf{Pattern:} a trite phrase or expression

\textbf{Use:} Use this for wording that has been repeated so often that it no longer sounds fresh.

\textbf{Example:} The introduction begins with a trite phrase about the importance of education.

\textbf{Pattern:} a trite remark or observation

\textbf{Use:} Use this for a comment that is familiar and lacks originality.

\textbf{Example:} The article ends with the trite observation that change is inevitable.

\textbf{Pattern:} sound or seem trite

\textbf{Use:} Use this when an idea may be sincere but its wording feels overused.

\textbf{Example:} It may sound trite, but careful preparation really does improve performance.

\textbf{Pattern:} avoid trite + noun

\textbf{Use:} Use this when recommending more original or precise language.

\textbf{Example:} Academic writers should avoid trite expressions when a more precise statement is possible.

\subsubsection{Common wrong usage}
\paragraph{Correction 1}
\textbf{Wrong:} The phrase is trite because it is easy to understand.

\textbf{Correct:} The phrase is trite because it has been used so often that it no longer sounds original.

\textbf{Why:} Trite does not mean simple or easy to understand. It means overused and lacking freshness.

\paragraph{Correction 2}
\textbf{Wrong:} The author made a tritely argument.

\textbf{Correct:} The author made a trite argument.

\textbf{Why:} Use the adjective trite before a noun such as argument.

\paragraph{Correction 3}
\textbf{Wrong:} This expression is very trite because nobody has used it before.

\textbf{Correct:} This expression is very original because nobody has used it before.

\textbf{Why:} A trite expression is familiar because it has been used many times.

\paragraph{Correction 4}
\textbf{Wrong:} The speaker trited several common ideas.

\textbf{Correct:} The speaker repeated several trite ideas.

\textbf{Why:} Trite is an adjective, not a verb.

\section{Word family}
\sectionrule
\subsection{triteness}
\textbf{Part of speech:} uncountable noun

\textbf{IPA:} /ˈtraɪtnəs/

\textbf{Definition:} the quality of being overused and no longer original or interesting

\textbf{Relationship to the headword:} This noun names the quality of being trite.

\textbf{Grammar pattern:} the triteness of + phrase, idea, remark, or theme

\textbf{Usage note:} It is much less common than trite.

\subsubsection{Examples}
\begin{enumerate}
  \item The triteness of the slogan made it easy to ignore.
  \item Critics complained about the triteness of the film's message.
\end{enumerate}

\subsection{tritely}
\textbf{Part of speech:} adverb

\textbf{IPA:} /ˈtraɪtli/

\textbf{Definition:} in an overused, predictable, or unoriginal way

\textbf{Relationship to the headword:} This adverb describes something said or expressed in a trite way.

\textbf{Grammar pattern:} tritely + say, conclude, observe, or describe

\textbf{Usage note:} It is uncommon; trite is much more useful for learners.

\subsubsection{Examples}
\begin{enumerate}
  \item The article tritely concludes that there are two sides to every issue.
  \item The character tritely describes love as the answer to every problem.
\end{enumerate}

\section{Common collocations}
\sectionrule
\subsection{trite phrase}
\textbf{Applies to:} Overused and no longer original

\textbf{Meaning:} a phrase that has been repeated so often that it has lost its freshness

\textbf{Example:} The essay begins with a trite phrase about living life to the fullest.

\subsection{trite expression}
\textbf{Applies to:} Overused and no longer original

\textbf{Meaning:} an overused way of expressing an idea

\textbf{Example:} The editor removed several trite expressions from the article.

\subsection{trite saying}
\textbf{Applies to:} Overused and no longer original

\textbf{Meaning:} a familiar saying that no longer seems fresh or insightful

\textbf{Example:} The speech was filled with trite sayings about hard work and success.

\subsection{trite remark}
\textbf{Applies to:} Overused and no longer original

\textbf{Meaning:} a predictable and unoriginal comment

\textbf{Example:} His trite remark added little to the discussion.

\subsection{trite observation}
\textbf{Applies to:} Overused and no longer original

\textbf{Meaning:} a familiar comment presented as an insight

\textbf{Example:} The author offers the trite observation that money cannot buy happiness.

\subsection{trite idea}
\textbf{Applies to:} Overused and no longer original

\textbf{Meaning:} an idea that has been repeated so often that it seems unoriginal

\textbf{Example:} A familiar topic does not have to produce a trite idea if it is developed carefully.

\subsection{trite argument}
\textbf{Applies to:} Overused and no longer original

\textbf{Meaning:} an argument that is familiar, predictable, and insufficiently original

\textbf{Example:} The article repeats the trite argument that all technological change is necessarily progress.

\subsection{sound trite}
\textbf{Applies to:} Overused and no longer original

\textbf{Meaning:} seem unoriginal because the wording is very familiar

\textbf{Pattern:} sound trite

\textbf{Example:} The advice may sound trite, but it is still useful.

\subsection{seem trite}
\textbf{Applies to:} Overused and no longer original

\textbf{Meaning:} appear predictable or overused

\textbf{Example:} The conclusion seems trite because similar statements appear throughout the literature.

\subsection{avoid trite expressions}
\textbf{Applies to:} Overused and no longer original

\textbf{Meaning:} choose not to use familiar expressions that have lost their impact

\textbf{Example:} Good writers avoid trite expressions and choose more precise language.

\section{Synonyms and antonyms}
\sectionrule
\subsection{Close synonyms}
\subsubsection{clichéd}
\textbf{Part of speech:} adjective

\textbf{Applies to:} Overused and no longer original

\textbf{Shared meaning:} Both describe something that has lost originality because it has been used too often.

\textbf{Difference:} Clichéd describes something that uses familiar clichés or predictable conventions. Trite is very close in meaning but can apply more broadly to remarks, ideas, observations, and advice that have become dull through repetition.

\textbf{Example:} The film was criticized for its clichéd characters and predictable ending.

\subsubsection{hackneyed}
\textbf{Part of speech:} adjective

\textbf{Applies to:} Overused and no longer original

\textbf{Shared meaning:} Both mean overused and no longer fresh or original.

\textbf{Difference:} Hackneyed is a stronger and somewhat more literary word for language or ideas that have become dull through excessive use. Trite is more common and straightforward.

\textbf{Example:} The writer relied on hackneyed images of heroic leaders and helpless victims.

\subsubsection{banal}
\textbf{Part of speech:} adjective

\textbf{Applies to:} Overused and no longer original

\textbf{Shared meaning:} Both can criticize ideas or remarks for lacking originality.

\textbf{Difference:} Banal emphasizes that an idea or remark is boring because it is ordinary and lacks originality. Trite more specifically emphasizes loss of freshness through repeated use.

\textbf{Example:} The interview produced several banal comments about the importance of teamwork.

\subsubsection{stale}
\textbf{Part of speech:} adjective

\textbf{Applies to:} Overused and no longer original

\textbf{Shared meaning:} Both can describe language or ideas that have lost their freshness.

\textbf{Difference:} Stale can figuratively describe ideas, jokes, or language that no longer feel fresh. Trite specifically suggests that repeated use has made them unoriginal and ineffective.

\textbf{Example:} The campaign relied on stale slogans that no longer attracted voters.

\subsection{Direct or useful antonyms}
\subsubsection{original}
\textbf{Part of speech:} adjective

\textbf{Applies to:} Overused and no longer original

\textbf{Opposition:} Original describes an idea or expression that is fresh or distinctive, while trite describes one that has lost its originality through repeated use.

\textbf{Example:} The student offered an original explanation that none of the other essays had considered.

\subsubsection{fresh}
\textbf{Part of speech:} adjective

\textbf{Applies to:} Overused and no longer original

\textbf{Opposition:} A fresh idea or expression feels new and interesting, while a trite one feels familiar and overused.

\textbf{Usage note:} This opposite applies to the figurative meaning of fresh.

\textbf{Example:} The author brings a fresh perspective to a familiar topic.

\section{Words learners may confuse}
\sectionrule
\subsection{cliché}
\textbf{Part of speech:} adjective versus countable noun

\textbf{Why learners confuse them:} The two words express almost the same basic idea but have different grammatical roles.

\textbf{Key difference:} Trite is an adjective describing something as overused and unoriginal. A cliché is a countable noun referring to an overused phrase, idea, or image itself.

\textbf{trite:} The conclusion sounds trite.

\textbf{cliché:} The phrase time heals all wounds is often treated as a cliché.

\subsection{trivial}
\textbf{Part of speech:} adjective

\textbf{Why learners confuse them:} Both are negative adjectives beginning with tri-, but they criticize different qualities.

\textbf{Key difference:} Trite means overused and unoriginal. Trivial means unimportant or too minor to deserve much attention.

\textbf{trite:} The speech contained several trite remarks about success.

\textbf{trivial:} The disagreement concerned a trivial detail and was resolved quickly.

\section{Etymology}
\sectionrule
\textbf{Origin:} Trite ultimately comes from Latin tritus, meaning rubbed or worn.

\textbf{Development:} The physical idea of something being worn down through repeated use developed into the modern figurative meaning of an expression or idea being worn out through repetition.

\textbf{Memory link:} A trite phrase is like a path that has been walked on so many times that it feels worn out.

\section{Practice exercises}
\sectionrule
\subsection{Word family choice}
Choose the best member of the word family for each blank.

\begin{enumerate}
  \item The \_\_\_\_\_ of the slogan made it ineffective despite its positive message.
  \item The article \_\_\_\_\_ concludes that people should simply learn from their mistakes.
  \item The speech contained several \_\_\_\_\_ observations about the importance of hard work.
\end{enumerate}

\subsection{Correct or incorrect?}
Decide whether each sentence is correct. Correct the inaccurate ones.

\begin{enumerate}
  \item The essay contains several trite phrases that have appeared in countless similar essays.
  \item The explanation is trite because it is simple and easy to understand.
  \item His advice sounded trite, although he sincerely believed it.
  \item The writer presented a tritely argument.
  \item The phrase is trite because nobody has ever used it before.
\end{enumerate}

\subsection{Collocation practice}
Complete each sentence with a natural collocation.

\begin{enumerate}
  \item The editor advised the writer to avoid trite \_\_\_\_\_ and use more precise language.
  \item The conclusion may \_\_\_\_\_ trite, but the underlying point is important.
  \item The author ends with the trite \_\_\_\_\_ that money cannot guarantee happiness.
  \item Students should avoid relying on \_\_\_\_\_ phrases when a more specific explanation is possible.
\end{enumerate}

\subsection{Complete answer key}
\subsubsection{Word family choice}
\begin{enumerate}
  \item \textbf{triteness.} The triteness of the slogan made it ineffective despite its positive message. \emph{Use the noun triteness for the quality of being overused and unoriginal.}
  \item \textbf{tritely.} The article tritely concludes that people should simply learn from their mistakes. \emph{Use the adverb tritely to describe how the conclusion is expressed.}
  \item \textbf{trite.} The speech contained several trite observations about the importance of hard work. \emph{Use the adjective trite before observations.}
\end{enumerate}

\subsubsection{Correct or incorrect?}
\begin{enumerate}
  \item \textbf{Correct} \emph{Trite correctly describes phrases that have become unoriginal through repeated use.}
  \item \textbf{Incorrect}. The explanation is simple and easy to understand. \emph{Trite does not mean simple. It means overused and no longer original or effective.}
  \item \textbf{Correct} \emph{An idea can be sincere or true and still sound trite because its wording is overused.}
  \item \textbf{Incorrect}. The writer presented a trite argument. \emph{Use the adjective trite before the noun argument.}
  \item \textbf{Incorrect}. The phrase is original because nobody has ever used it before. \emph{Trite describes something that has been used too often, not something completely new.}
\end{enumerate}

\subsubsection{Collocation practice}
\begin{enumerate}
  \item \textbf{expressions.} The editor advised the writer to avoid trite expressions and use more precise language. \emph{Trite expressions are familiar expressions that have lost their freshness through repeated use.}
  \item \textbf{sound.} The conclusion may sound trite, but the underlying point is important. \emph{Sound trite means seem unoriginal or overused.}
  \item \textbf{observation.} The author ends with the trite observation that money cannot guarantee happiness. \emph{A trite observation is a familiar comment presented as an insight.}
  \item \textbf{trite.} Students should avoid relying on trite phrases when a more specific explanation is possible. \emph{A trite phrase has been repeated so often that it no longer sounds fresh or original.}
\end{enumerate}

\vfill
\begin{center}
\small\color{Muted} Created and compiled by \textbf{Amir Kooshky}\\
\href{https://t.me/KooshkyTOEFL}{Telegram Channel}\quad\textbullet\quad
\href{https://www.instagram.com/kooshkytoefl}{Instagram}\quad\textbullet\quad
\href{https://t.me/KooshkyTOEFL\_pv}{Personal Telegram}
\end{center}
\end{document}
